\chapter{Neutron Scattering II: Partial Waves}

\sourcepage{1}
\subsection*{From last time}
\begin{enumerate}
  \item Reactions, Bohr, and the compound nucleus.
  \item Cross sections.
  \item $\widehat H\psi=E\psi$ for a two-body system and a central potential $V(r)$.
\end{enumerate}
The probability current and continuity equation were left as
\[
  \bm j=\frac{\hbar}{2m\ii}\left(\overline\psi\nabla\psi-\psi\nabla\overline\psi\right),
  \qquad
  \frac{\partial}{\partial t}(\overline\psi\psi)+\nabla\cdot\bm j=0.
\]
The topics for the present notes are quantum-mechanical scattering and $\sigma(\theta)$, partial waves for $V(\bm r)=V(r)$, and a multichannel overview.

\sourcepage{2}
The differential cross section is
\[
  \sigma(\Omega)=\frac{\dd\sigma}{\dd\Omega}
  =\frac{\text{number scattered}}
  {\text{target}\times\text{incident flux}},
\]
with units of barns per steradian.  For
\[
  \psi_{\mathrm{in}}=A\ee^{\ii kz},
  \qquad
  \psi_{\mathrm{sc}}=A f(\Omega)\frac{\ee^{\ii kr}}{r},
\]
the currents are
\[
  j_i=v_i|A|^2,
  \qquad
  j_f=v_f|A|^2\frac{|f|^2}{R^2}.
\]
Particle balance through an element of solid angle gives
\[
  j_i\frac{\dd\sigma}{\dd\Omega}=j_fR^2,
\]
and hence
\begin{equation}
  \frac{\dd\sigma}{\dd\Omega}=\frac{v_f}{v_i}|f|^2.
  \label{eq:ns2-diff-xs}
\end{equation}

\sourcepage{3}
From the previous discussion, a plane wave incident along $z$ and scattered by a spherical potential has the asymptotic form
\[
  \psi_{\mathrm{asym}}\sim \ee^{\ii kz}+f(\theta,\phi)\frac{\ee^{\ii kr}}{r}.
\]
The incident plane wave has cylindrical symmetry about $z$, and the spherical potential cannot break that symmetry.  Therefore
\[
  f(\theta,\phi)=f(\theta),
  \qquad
  \sigma(\theta,\phi)=\sigma(\theta).
\]
Given a potential $V(r)$, the Schr\"odinger equation can in principle be solved to obtain the cross section.  Two general schemes are noted:
\begin{enumerate}
  \item the method of partial waves, an appropriate expansion into an infinite set of ordinary differential equations;
  \item the Born series, an infinite sum whose first term gives the Born approximation.
\end{enumerate}
The notes remark that this is a common dichotomy in mathematical physics.

\sourcepage{4}
\section{Method of partial waves}
The radial equation is recalled in the form
\begin{equation}
  \left[-\frac{\hbar^2}{2m}
  \left(\frac{\dd^2}{\dd r^2}+\frac{2}{r}\frac{\dd}{\dd r}
  -\frac{\ell(\ell+1)}{r^2}\right)
  +V(r)-E\right]\psi_\ell(r)=0.
  \label{eq:ns2-radial-psi}
\end{equation}
Expand the axially symmetric wave function in Legendre polynomials,
\begin{equation}
  \psi(r,\theta)
  =\sum_{\ell=0}^{\infty}(2\ell+1)\ii^\ell
  P_\ell(\cos\theta)\frac{1}{kr}\chi_\ell(r).
  \label{eq:ns2-pw-expansion}
\end{equation}
The factor $1/(kr)$ is included for convenience.  For $V(r)=0$, $\chi_\ell$ is related to the spherical Bessel and Neumann functions $j_\ell$ and $y_\ell$.

\sourcepage{5}
Substitution of \eqref{eq:ns2-pw-expansion} into the Schr\"odinger equation, followed by multiplication by $P_m(\cos\theta)$ and integration over $\mu=\cos\theta$, uses
\[
  \int_{-1}^{1}P_m(\mu)P_\ell(\mu)\dd\mu
  =\frac{2}{2\ell+1}\delta_{m\ell}.
\]
The result is one independent radial equation for each partial wave:
\begin{equation}
  \left[-\frac{\hbar^2}{2m}
  \left(\frac{\dd^2}{\dd r^2}-\frac{\ell(\ell+1)}{r^2}\right)
  +V(r)-E\right]\chi_\ell(r)=0.
  \label{eq:ns2-chi-equation}
\end{equation}

\sourcepage{6}
Each partial-wave equation requires two boundary conditions:
\begin{enumerate}
  \item For a finite wave function,
  \[
    \chi_\ell(0)=0.
  \]
  \item At some radius $r=a$ outside the range of the potential, match $\chi_\ell$ to its asymptotic form.  This assumes a short-range potential, $V(r)\simeq0$ for $r>a$.
\end{enumerate}
In vacuum, the solutions are proportional to $r h_\ell^{\pm}(kr)$.  The plane wave must be expressed in the same basis.  Since
\[
  \ee^{\ii kz}=\ee^{\ii kr\cos\theta}=\ee^{\ii\rho\cos\theta},
\]
write
\[
  \ee^{\ii\rho\cos\theta}
  =\sum_{\ell=0}^{\infty}(2\ell+1)a_\ell(\rho)P_\ell(\cos\theta).
\]
Orthogonality gives
\[
  \frac{2}{2\ell+1}a_\ell(\rho)
  =\int_{-1}^{1}\ee^{\ii\rho\mu}P_\ell(\mu)\dd\mu
  =2\ii^\ell j_\ell(\rho),
\]
so
\begin{equation}
  \ee^{\ii kz}
  =\sum_{\ell=0}^{\infty}(2\ell+1)\ii^\ell
  j_\ell(kr)P_\ell(\cos\theta).
  \label{eq:ns2-plane-wave-expansion}
\end{equation}

\sourcepage{7}
Equivalently,
\[
  \ee^{\ii\rho\cos\theta}
  =\sum_{\ell=0}^{\infty}(2\ell+1)a_\ell(\rho)P_\ell(\cos\theta),
  \qquad
  a_\ell(\rho)=\frac12\int_{-1}^{1}P_\ell(\mu)\ee^{\ii\rho\mu}\dd\mu.
\]
Using the integral representation
\[
  j_\ell(\rho)=\frac{(-\ii)^\ell}{2}
  \int_{-1}^{1}\ee^{\ii\rho\mu}P_\ell(\mu)\dd\mu,
\]
one again finds $a_\ell=\ii^\ell j_\ell$ and obtains \eqref{eq:ns2-plane-wave-expansion}.

\sourcepage{8}
Since
\[
  h_\ell^\pm=j_\ell\pm\ii y_\ell,
  \qquad
  j_\ell=\frac12(h_\ell^-+h_\ell^+),
\]
the incident plane wave may be written
\[
  \ee^{\ii kz}
  =\frac12\sum_{\ell=0}^{\infty}(2\ell+1)\ii^\ell
  \left[h_\ell^-(kr)+h_\ell^+(kr)\right]P_\ell(\cos\theta).
\]
Thus the incident beam consists of incoming and outgoing radial waves of equal amplitude: it is a plane wave that approaches the target and leaves it unchanged.

Outside the potential, the asymptotic partial-wave solution must be a linear combination of the two basic solutions.  The notes write
\begin{equation}
  \chi_\ell^{\mathrm{ext}}(r)
  =A_\ell\left[H_\ell^-(kr)-S_\ell H_\ell^+(kr)\right],
  \label{eq:ns2-exterior-solution}
\end{equation}
where $A_\ell$ and $S_\ell$ are complex coefficients.  The coefficient $S_\ell$ is the $S$-matrix element for the $\ell$th partial wave.
% The source parenthetical appears to read ``S_l is zero for V=0,'' although the subsequent equations use S_l=1 for no scattering.

\sourcepage{9}
Within the potential, solve for a radial wave function $u_\ell(r)$.  Its overall normalization is arbitrary, so the actual solution is
\[
  \chi_\ell(r)=B u_\ell(r).
\]
At the matching radius $a$,
\[
  \chi_\ell(a)=\chi_\ell^{\mathrm{ext}}(a),
  \qquad
  \chi_\ell'(a)=\chi_\ell^{\mathrm{ext}\prime}(a).
\]
The notes define the dimensionless logarithmic-derivative quantity
\begin{equation}
  R_\ell
  =\frac{1}{a}\frac{\chi_\ell(a)}{\chi_\ell'(a)}
  =\frac{1}{a}\frac{u_\ell(a)}{u_\ell'(a)}.
  \label{eq:ns2-R-def}
\end{equation}
This is the $R$-matrix element of the $\ell$th partial wave.  Combining \eqref{eq:ns2-exterior-solution} and \eqref{eq:ns2-R-def},
\[
  R_\ell
  =\frac{1}{a}
  \frac{H_\ell^-(ka)-S_\ell H_\ell^+(ka)}
  {H_\ell^{-\prime}(ka)-S_\ell H_\ell^{+\prime}(ka)}.
\]
Solving for $S_\ell$ gives
\begin{equation}
  S_\ell
  =\frac{H_\ell^-(ka)-aR_\ell H_\ell^{-\prime}(ka)}
  {H_\ell^+(ka)-aR_\ell H_\ell^{+\prime}(ka)}.
  \label{eq:ns2-S-from-R}
\end{equation}
The $S$-matrix element is independent of the normalization of the interior solution and is uniquely defined by the potential.

\sourcepage{10}
To compute the scattering amplitude, use the asymptotic partial-wave form
\[
  \psi^{\mathrm{asym}}
  \sim \frac{1}{kr}\sum_{\ell=0}^{\infty}(2\ell+1)\ii^\ell
  P_\ell(\cos\theta)A_\ell
  \left[H_\ell^-(kr)-S_\ell H_\ell^+(kr)\right].
\]
For $kr\gg1$,
\[
  H_\ell^\pm(kr)\sim \ii^{\mp(\ell+1)}\ee^{\pm\ii kr}.
\]
Compare with
\[
  \psi^{\mathrm{asym}}=\ee^{\ii kz}+f(\theta)\frac{\ee^{\ii kr}}{r}
\]
and match the $\ee^{-\ii kr}$ and $\ee^{\ii kr}$ terms.  Matching the incoming term gives
\[
  A_\ell=\frac{\ii}{2}.
\]

\sourcepage{11}
Expand the amplitude as
\[
  f(\theta)=\sum_{\ell=0}^{\infty}(2\ell+1)f_\ell P_\ell(\cos\theta),
  \qquad
  f_\ell=\frac12\int_{-1}^{1}P_\ell(\mu)f(\mu)\dd\mu.
\]
Matching outgoing terms and using orthogonality yields
\begin{equation}
  f_\ell=\frac{1}{2k\ii}(S_\ell-1).
  \label{eq:ns2-fl}
\end{equation}
Therefore
\begin{equation}
  f(\theta)=\frac{1}{2k\ii}
  \sum_{\ell=0}^{\infty}(2\ell+1)P_\ell(\cos\theta)(S_\ell-1),
  \label{eq:ns2-f-S}
\end{equation}
and the differential cross section is
\begin{equation}
  \sigma(\theta)
  =\frac{1}{4k^2}
  \left|\sum_{\ell=0}^{\infty}(2\ell+1)P_\ell(\cos\theta)(S_\ell-1)\right|^2.
  \label{eq:ns2-diff-xs-S}
\end{equation}

\sourcepage{12}
The asymptotic partial waves are linear combinations of $H_\ell^+$ and $H_\ell^-$.  A more physical parametrization is
\begin{equation}
  S_\ell=\ee^{2\ii\delta_\ell},
  \qquad
  \delta_\ell=\frac{1}{2\ii}\ln S_\ell+n\pi,
  \quad n\in\mathbb Z.
  \label{eq:ns2-phase-shift}
\end{equation}
Then
\begin{equation}
  f(\theta)=\frac{1}{k}
  \sum_{\ell=0}^{\infty}(2\ell+1)P_\ell(\cos\theta)
  \ee^{\ii\delta_\ell}\sin\delta_\ell.
  \label{eq:ns2-f-phase}
\end{equation}
The exterior partial wave can likewise be written
\[
  \chi_\ell^{\mathrm{ext}}
  =\ee^{\ii\delta_\ell}
  \left[\cos\delta_\ell\,F_\ell(kr)
  +\sin\delta_\ell\,G_\ell(kr)\right]
  \sim \ee^{\ii\delta_\ell}
  \sin\left(kr+\delta_\ell-\frac{\ell\pi}{2}\right).
\]
Hence the name \emph{phase shift}.  The sketch notes that $\delta_\ell>0$ corresponds to a weakly attractive potential.

\sourcepage{13}
The phase-shift representation gives several useful results.  The elastic cross section for one partial wave is
\begin{equation}
  \sigma_{\mathrm{el},\ell}
  =\pi\bar\lambda^2(2\ell+1)|1-S_\ell|^2,
  \qquad \bar\lambda=\frac{1}{k}.
  \label{eq:ns2-el-partial}
\end{equation}
The total elastic cross section is
\begin{equation}
  \sigma_{\mathrm{el}}
  =\pi\bar\lambda^2
  \sum_{\ell=0}^{\infty}(2\ell+1)|1-S_\ell|^2.
  \label{eq:ns2-el-total}
\end{equation}
The forward amplitude obeys the optical theorem in the form
\begin{equation}
  \operatorname{Im}f(0)=\frac{k}{4\pi}\sigma_{\mathrm{el}}
  \label{eq:ns2-optical-elastic}
\end{equation}
for purely elastic scattering.

Since the exterior partial wave is proportional to $H_\ell^- -S_\ell H_\ell^+$, the difference between incoming and outgoing radial intensities is $1-|S_\ell|^2$.  Define the reaction cross section
\begin{equation}
  \sigma_R=\frac{\pi}{k^2}
  \sum_{\ell=0}^{\infty}(2\ell+1)(1-|S_\ell|^2).
  \label{eq:ns2-reaction-xs}
\end{equation}
It vanishes when $|S_\ell|=1$.

\sourcepage{14}
\subsection*{Net current through a large sphere}
The probability current is
\[
  \bm j=\frac{\hbar}{2m\ii}
  (\overline\psi\nabla\psi-\psi\nabla\overline\psi).
\]
At a large radius $R$, the net radial current is obtained from $\widehat{\bm r}\cdot\bm j$ and integrated over the sphere.  Using the asymptotic partial-wave form and Legendre orthogonality, the $\ell$th-wave contribution is proportional to the outgoing-minus-incoming intensity,
\[
  \left(1-|S_\ell|^2\right).
\]
The notes reduce the angular integral with
\[
  \int_0^{2\pi}\dd\phi\int_{-1}^{1}P_\ell^2(\mu)\dd\mu
  =\frac{4\pi}{2\ell+1}.
\]

\sourcepage{15}
The total net outward current is recorded as
\[
  j_{\mathrm{out}}
  =-\frac{\hbar\pi}{mk}
  \sum_{\ell=0}^{\infty}(2\ell+1)(1-|S_\ell|^2).
\]
The reaction cross section is defined as the loss of current divided by the incident current.  For $\psi_{\mathrm{in}}=\ee^{\ii kz}$,
\[
  j_{\mathrm{in}}=v=\frac{\hbar k}{m},
\]
which recovers
\[
  \sigma_R=-\frac{j_{\mathrm{out}}}{j_{\mathrm{in}}}
  =\pi\bar\lambda^2
  \sum_{\ell=0}^{\infty}(2\ell+1)(1-|S_\ell|^2).
\]
Thus
\begin{align}
  \sigma_{R,\ell}
  &=\pi\bar\lambda^2(2\ell+1)(1-|S_\ell|^2),\\
  \sigma_R
  &=\pi\bar\lambda^2\sum_{\ell=0}^{\infty}(2\ell+1)(1-|S_\ell|^2).
\end{align}
The total cross section is
\begin{align}
  \sigma_{T,\ell}
  &=\sigma_{\mathrm{el},\ell}+\sigma_{R,\ell}\\
  &=\pi\bar\lambda^2(2\ell+1)
  \left(|1-S_\ell|^2+1-|S_\ell|^2\right)\\
  &=2\pi\bar\lambda^2(2\ell+1)
  \left[1-\operatorname{Re}(S_\ell)\right],\\
  \sigma_T
  &=2\pi\bar\lambda^2\sum_{\ell=0}^{\infty}(2\ell+1)
  \left[1-\operatorname{Re}(S_\ell)\right].
\end{align}

\sourcepage{16}
Return to \eqref{eq:ns2-S-from-R}.  Dropping the partial-wave subscript temporarily,
\[
  S
  =\frac{1-a(H^{-\prime}/H^-)R}{1-a(H^{+\prime}/H^+)R}
  \frac{H^-}{H^+}.
\]
Writing $H^+=G+\ii F$,
\begin{align*}
  a\frac{H^{+\prime}}{H^+}
  &=a\frac{G'+\ii F'}{G+\ii F}\\
  &=a\frac{GG'+FF'}{G^2+F^2}
  +\ii a\frac{GF'-FG'}{G^2+F^2}\\
  &\equiv \widehat S+\ii\widehat P,
\end{align*}
where $\widehat S$ is the shift function and $\widehat P$ the penetrability.  For the first partial waves, the notes list
\begin{center}
\begin{tabular}{c c c}
\toprule
$\ell$ & $\widehat S_\ell(\rho)$ & $\widehat P_\ell(\rho)$\\
\midrule
$0$ & $0$ & $\rho$\\
$1$ & $-1/(\rho^2+1)$ & $\rho^3/(\rho^2+1)$\\
\bottomrule
\end{tabular}
\end{center}
Also define the hard-sphere phase $\Delta$ by
\[
  \frac{H^-}{H^+}=\frac{G-\ii F}{G+\ii F}=\ee^{2\ii\Delta}.
\]
The resulting $S$-matrix expression is recorded in the notes as a ratio involving $1-R\widehat S$ and $R\widehat P$, multiplied by $\ee^{2\ii\Delta}$.  Consistent with the immediately following phase decomposition, one may write
\begin{equation}
  S_\ell
  =\frac{1-R_\ell\widehat S_\ell+\ii R_\ell\widehat P_\ell}
  {1-R_\ell\widehat S_\ell-\ii R_\ell\widehat P_\ell}
  \ee^{2\ii\Delta_\ell}.
  \label{eq:ns2-S-shift-penetration}
\end{equation}
\editorial{The handwritten signs in the boxed fraction are difficult to distinguish; the displayed form follows the algebra and the phase relation on the next source page.}

\sourcepage{17}
For elastic scattering,
\[
  \sigma_{\mathrm{el},\ell}
  =\pi\bar\lambda^2(2\ell+1)|1-S_\ell|^2.
\]
Let
\[
  1-R_\ell\widehat S_\ell=\cos\alpha_\ell,
  \qquad
  R_\ell\widehat P_\ell=\sin\alpha_\ell,
\]
up to an arbitrary common normalization.  Then
\[
  \frac{\cos\alpha_\ell+\ii\sin\alpha_\ell}
  {\cos\alpha_\ell-\ii\sin\alpha_\ell}
  =\frac{1+\ii\tan\alpha_\ell}{1-\ii\tan\alpha_\ell}
  =\ee^{2\ii\alpha_\ell}.
\]
Consequently,
\begin{equation}
  S_\ell=\ee^{2\ii\alpha_\ell+2\ii\Delta_\ell},
  \qquad
  \delta_\ell=\alpha_\ell+\Delta_\ell.
  \label{eq:ns2-compound-hard-sphere-phase}
\end{equation}
The notes label $\alpha_\ell$ as the compound-elastic contribution and $\Delta_\ell$ as the hard-sphere or potential contribution.

\subsection*{Exercise}
Write $1-S_\ell$ as $\ee^{2\ii\Delta_\ell}$ times a sum of potential and compound-elastic amplitudes.  Then show or investigate:
\begin{enumerate}
  \item the decomposition of $\sigma_{\mathrm{el},\ell}$ into potential-elastic, compound-elastic, and interference contributions;
  \item the total partial-wave cross section as a function of the corresponding amplitudes;
  \item for $\ell=0$, the low-energy behavior $\sigma_{T,0}\propto 1/v$ in the absorptive case indicated in the notes.
\end{enumerate}
